% page 581 du fac-simile (image p-580 du scan)
% bloc 170.2 x 267.9 mm ; marge gauche 31.8 mm
\pgc{10.160mm}{0.000mm}{
\l{}
\l{}
\l{}
\l{}
\l{\cel{56}573}
\l{}
\l{\sou{telegonio}. (\sou{biol.}) Transmiso, a la genituro ek mariajo dato-}
\l{\cel{3}duesma, di karakteri qui devenas de la patrulo unesma-mariaja.}
\l{\cel{3}- DEFIS.}
\l{}
\l{\sou{telegrafar}. (\sou{trans., ad)}. Sendar (texto) per aparato qua transmi-}
\l{\cel{3}sas la depeshi preske instantale. - DEFISR.}
\l{}
\l{\sou{telegramo}.\cel{2}Depesho tranmisita per telegrafo. - DEFIRS.}
\l{}
\l{\sou{teleologio}. (\sou{biol.)}\cel{2}Doktrino segun qua la enti esas organizita}
\l{\cel{3}por la skopo a qua li destinesis. - DEFIRS.}
\l{}
\l{\sou{telepatio}. (\sou{metapsik.}) Nomo per qua onu indikas ula senti (sen-}
\l{\cel{3}sala o psikala) de persono, pri ulo reala quon agis sama-tempe}
\l{\cel{3}altra persono, ma ye tanta disto ed en cirkonstanci tala, ke}
\l{\cel{3}olua saveso da la persono unesme-mencionita aspektas neposibla.}
\l{\cel{3}- DEFIRS.}
\l{}
\l{\sou{teleskopo}. Optik-instrumento, uzata por observar la objekti tre}
\l{\cel{3}fora, per la reflektigo, sur speguli, di la imajo de ta objekti,}
\l{\cel{3}e per plugrandigar ta reflekturo per lensi. - DEFIRS.}
\l{}
\l{telofazo. (\sou{biol.}) Fazo finala di la mitoso. -}
\l{}
\l{\sou{telolecito}. (\sou{biol.)}\cel{2}Dicesas pri la ovulo kun segmentesko partala}
\l{\cel{3}(uceli, repteri, fishi) en qua la vitelo formaciva esas aparta}
\l{\cel{3}de la vitelo nutriva.}
\l{}
\l{\sou{teluro}. (\sou{kemio)}.\cel{2}Korpo simpla, metala, deskovrita en 1782 en}
\l{\cel{3}la oro-mineyi di Transilvania. -\cel{1}\sou{Simbolo kemiala}\cel{1}: Te.}
\l{}
\l{\sou{temo}. Objekto di studio, propoziciono quan onu traktas, develo-}
\l{\cel{3}pas. - DEFIRS.}
\l{}
\l{\sou{tempo.}- I.Duro di to quo havas komenco-punto e fino-parto.}
\l{\cel{3}II. (\sou{gram.)}\cel{1}Irga ek la flexioni diversa di un modo, qua indi-}
\l{\cel{3}kas ke la ago o la subiso quan expresas la verbo esas prezenta,}
\l{\cel{3}pasinta o futura. - III. (\sou{muziko}).Unajo dividala di la mezuro}
\l{\cel{3}muzikala. - EFIS.}
\l{}
\l{\sou{tema}to. (\sou{muziko)}. Temo, motivo, melodiala di muzikajo. - II.}
\l{\cel{3}(\sou{gram.}) Parto nevariiva di vorto, qua recevas la flexioni.}
\l{\cel{3}- DEFIRS.}
\l{}
\l{\sou{temerara}. Di qua la audaco, la sen-timeso, esas ne-reflektita.EFIS}
\l{}
\l{\sou{temperar}. (\sou{trans.}) Moderar per ta o ca pludolcigo, diminuto ;}
\l{\cel{3}igar min strikta, min rigoroza. - DEFIS.}
\l{}
\l{temperamento. (\sou{fiziol}.) Konstituciono propra ye korpo organizita,}
\l{\cel{3}qua dependas de la maniero ye qua la funcioni inter-equilibr-}
\l{\cel{3}esas en lu. - DEFIRS.}
\l{}
\l{\cel{1}\sou{temperaturo}. Varmeso-grado di korpo, di loko, evaluebla per la}
\l{\cel{4}senti quin lu produktas ye nia organi, o lua efiko a termo-}
\l{\cel{4}metro. - DEFIRS.}
}
